\documentclass[11pt]{article}
\usepackage[margin=1in]{geometry}
\usepackage{amsmath}
\usepackage{amssymb}
\usepackage{graphicx}
\usepackage{booktabs}
\usepackage{xcolor}
\usepackage{hyperref}

\title{ConstraintServe: Optimal Model Selection and Resource Allocation for Multi-Model LLM Inference}
\author{Kirti Chavhan \\ Independent Researcher \\ \texttt{kirti.a.chavhan@gmail.com}}
\date{September 2026}

\begin{document}

\maketitle

\begin{abstract}
This paper addresses the challenge of serving multiple heterogeneous Large Language Models (LLMs) efficiently on shared hardware infrastructure. Production deployments require serving dozens of different models with varying resource requirements, latencies, costs, and specialized capabilities, yet existing systems either optimize for single models or use heuristic-based approaches. We propose \textbf{ConstraintServe}, an advanced system combining optimal model selection with dynamic resource allocation and fairness-aware scheduling. The system handles 10 per-request parameters including task type, token count, context window, and specialized capabilities. Through discrete-event simulation with five independent runs, we demonstrate \textbf{88.4\% cost reduction} (95\% CI: 87.2\%-89.6\%) compared to baseline approaches, achieving \textbf{95.41\% SLA compliance} (95\% CI: 94.82\%-96.00\%) across diverse priority levels and tasks. Task-aware routing shows specialized optimization: code generation costs \$97.54, while search queries cost only \$10.00 per 1000 requests. The constraint satisfaction approach demonstrates robustness across varying workloads, cost models, quality requirements, and token lengths.
\end{abstract}

\section{Introduction}

\subsection{The Multi-Model Serving Challenge}

Modern AI companies operate dozens to hundreds of different Large Language Models in production simultaneously, with significant variations:

\begin{itemize}
    \item \textbf{Model size:} 7B to 70B parameters (10x memory variance)
    \item \textbf{Latency spectrum:} 48ms to 310ms per inference (6.5x range)
    \item \textbf{Cost variance:} Estimated \$0.01 to \$0.50 per request (50x difference)
    \item \textbf{Accuracy diversity:} Specialized models outperform general models on specific tasks
    \item \textbf{Capability heterogeneity:} Vision, tool use, function calling vary by model
\end{itemize}

\textbf{The Fundamental Problem:} Given N models with different capabilities and limited GPU hardware (typically 40-80GB memory), how do you select which model to process each incoming request, considering not just quality, latency, and cost, but also task type, token length, required capabilities, and session consistency?

\subsection{Why Existing Systems Are Inadequate}

\begin{table}[h]
\centering
\begin{tabular}{llll}
\toprule
\textbf{System} & \textbf{Strength} & \textbf{Limitation} \\
\midrule
vLLM & Excellent single-model optimization & Cannot select between models \\
Ray Serve & Distributed multi-model support & Heuristic-based routing; no task awareness \\
Clipper & Model selection exists & Offline only; ignores specialized capabilities \\
ConstraintServe & Task/capability-aware + optimal routing & Novel contribution \\
\bottomrule
\end{tabular}
\end{table}

\subsection{Our Contribution}

We propose \textbf{ConstraintServe}, a three-component system with enhanced parameter handling:

\begin{enumerate}
    \item \textbf{Task-Aware Model Selection} -- Real-time optimal model choice considering task type and specialized capabilities
    \item \textbf{Dynamic Resource Allocation} -- GPU memory allocation under hard constraints
    \item \textbf{Fairness Scheduling} -- SLA-aware priority-based request scheduling with session consistency
\end{enumerate}

\textbf{Key Innovation:} First system to combine constraint satisfaction with task-aware routing, token-aware latency prediction, and capability-aware filtering for multi-model LLM serving.

\section{Related Work}

\subsection{LLM Serving Infrastructure}

vLLM (Kwon et al., 2023) revolutionized single-model LLM serving through PagedAttention memory optimization. ConstraintServe builds on vLLM's efficient infrastructure while extending it to heterogeneous multi-model scenarios with task awareness.

Ray Serve (Nishihara et al., 2020) enables distributed model serving with flexible routing policies. However, routing decisions use heuristics rather than optimization, and do not consider task type or specialized capabilities.

\subsection{Model Selection}

Clipper (Crankshaw et al., 2017) introduced adaptive model selection using offline ensemble learning. Key differences from ConstraintServe:
\begin{itemize}
    \item Clipper trains on historical data; ConstraintServe makes real-time decisions based on per-request constraints
    \item Clipper ignores task type; ConstraintServe uses task-aware routing with affinity scores
    \item Clipper doesn't handle specialized capabilities (vision, tools, function calling)
\end{itemize}

\subsection{Research Gap}

\textbf{No prior system combines:} Task-aware real-time model selection + token-aware latency prediction + capability filtering + dynamic resource allocation + SLA-aware fairness scheduling for multi-model LLM inference.

\section{Problem Formulation}

We formalize multi-model LLM serving as a constraint satisfaction problem with 10 parameters.

\subsection{Formal Definition}

\textbf{Given:}
\begin{itemize}
    \item $M = \{m_1, m_2, \ldots, m_n\}$: Set of available models with properties:
    \begin{itemize}
        \item accuracy, latency, cost, memory
        \item task support, context window
        \item capability flags (vision, tools, function calling)
    \end{itemize}
    \item $K$: Total GPU memory available (GB)
    \item Request stream $R$ with 10 per-request parameters
\end{itemize}

\textbf{Objective:} For each request $r$, select model $m^*$ to:

\begin{equation}
\text{Minimize: } \text{cost}(m^*) \times (1 + \text{latency\_penalty}(m^*)) / (1 + \text{task\_affinity}(m^*, r))
\end{equation}

Subject to hard constraints:
\begin{align}
\text{accuracy}(m^*) &\geq \text{quality\_min}(r) \\
\text{latency}(m^*, r.\text{tokens}) &\leq \text{latency\_sla}(r) \\
\text{cost}(m^*) &\leq \text{cost\_budget}(r) \\
m^*.\text{task} &\in \text{supported\_tasks}(m^*) \\
m^*.\text{context} &\geq r.\text{min\_context\_window} \\
m^*.\text{supports\_vision} &\geq r.\text{needs\_vision} \\
m^*.\text{supports\_tools} &\geq r.\text{needs\_tools} \\
m^*.\text{supports\_fc} &\geq r.\text{needs\_function\_calling} \\
\text{total\_allocated\_memory} &\leq K \\
m^* &\in M
\end{align}

\subsection{Request Parameters (10 Total)}

\textbf{Tier 1 - Core Task Parameters:}
\begin{enumerate}
    \item \textbf{task\_type}: ``code'', ``chat'', ``translate'', ``summarize'', ``search''
    \item \textbf{expected\_tokens}: Integer (50--1000 typical)
    \item \textbf{consistency\_required}: Boolean (session consistency)
    \item \textbf{min\_context\_window}: Integer (4k, 8k, 32k, 128k tokens)
\end{enumerate}

\textbf{Tier 2 - Capability Requirements:}
\begin{enumerate}
    \setcounter{enumi}{4}
    \item \textbf{needs\_vision}: Boolean
    \item \textbf{needs\_tools}: Boolean
    \item \textbf{needs\_function\_calling}: Boolean
    \item \textbf{preferred\_batch\_size}: Integer (1, 8, 32)
\end{enumerate}

\textbf{Plus:}
\begin{enumerate}
    \setcounter{enumi}{8}
    \item \textbf{timeout\_tolerance}: ``hard'' or ``soft''
    \item \textbf{multimodal\_input}: ``text'', ``image'', ``audio''
\end{enumerate}

\section{System Design}

\subsection{Enhanced Model Selection Algorithm}

The algorithm operates in O(n) time where $n \leq 50$ models:

\begin{enumerate}
    \item \textbf{Session Check:} If session\_id has sticky assignment, return that model
    \item \textbf{Filter:} Identify feasible models meeting ALL hard constraints:
    \begin{itemize}
        \item Accuracy $\geq$ quality\_min
        \item Latency (token-aware) $\leq$ latency\_sla
        \item Cost $\leq$ budget
        \item GPU memory available
        \item Task support includes requested task
        \item Context window $\geq$ min\_context\_window
        \item All capability requirements met
    \end{itemize}
    \item \textbf{Score:} Rank feasible models by:
    $$\text{score} = \frac{\text{cost} \times (1 + \text{latency}/500)}{1 + \text{task\_affinity}(m, \text{task})}$$
    \item \textbf{Select:} Choose minimum-score model
    \item \textbf{Session Sticky:} If session\_id and consistency\_required, record assignment
\end{enumerate}

Decision time: $<0.01$ms on typical hardware.

\subsection{Task-Aware Routing}

Each model has task-specific affinity scores. Example:
\begin{table}[h]
\centering
\begin{tabular}{llllll}
\toprule
\textbf{Model} & \textbf{Code} & \textbf{Chat} & \textbf{Translate} & \textbf{Summarize} & \textbf{Search} \\
\midrule
Llama-2-7B & 0.65 & 0.90 & 0.82 & 0.85 & 0.87 \\
Mistral-7B & 0.75 & 0.92 & 0.88 & 0.89 & 0.88 \\
Llama-2-13B & 0.85 & 0.94 & 0.93 & 0.92 & 0.91 \\
Code-Llama-34B & 0.98 & 0.85 & 0.80 & 0.83 & 0.82 \\
Llama-2-70B & 0.95 & 0.96 & 0.97 & 0.95 & 0.93 \\
\bottomrule
\end{tabular}
\end{table}

Task affinity directly influences model selection: code generation prefers Code-Llama-34B, while search queries prefer small fast models (Llama-2-7B).

\subsection{Token-Aware Latency Prediction}

Model latency depends on token count:

$$\text{latency}(m, \text{tokens}) = \text{base\_latency}(m) \times (1 + 0.05 \times (\text{tokens}/100 - 1))$$

This enables more accurate selection for requests with varying token lengths.

\section{Experimental Methodology}

\subsection{Simulation Framework}

We implement a discrete-event simulator that:
\begin{enumerate}
    \item Simulates request arrivals (Poisson process, 100ms avg inter-arrival)
    \item Extracts 10 parameters from each request
    \item Runs enhanced model selection algorithm
    \item Simulates GPU resource allocation with session tracking
    \item Records detailed metrics per task type
\end{enumerate}

\subsection{Setup}

\textbf{Simulated GPU Hardware:} 40GB (A100-equivalent)

\textbf{Model Zoo with Capabilities:}

\begin{table}[h]
\centering
\begin{tabular}{lrrrrccc}
\toprule
\textbf{Model} & \textbf{Mem} & \textbf{Lat} & \textbf{Acc} & \textbf{Cost} & \textbf{Vision} & \textbf{Tools} & \textbf{FC} \\
\midrule
Llama-2-7B & 14GB & 52ms & 0.88 & \$0.01 & No & Yes & No \\
Mistral-7B & 14GB & 48ms & 0.87 & \$0.02 & No & Yes & Yes \\
Llama-2-13B & 26GB & 95ms & 0.92 & \$0.05 & No & Yes & Yes \\
Code-Llama-34B & 28GB & 180ms & 0.94 & \$0.15 & No & Yes & Yes \\
Llama-2-70B & 32GB & 310ms & 0.96 & \$0.50 & Yes & Yes & Yes \\
\bottomrule
\end{tabular}
\end{table}

\textbf{Workload:} 10,000 synthetic requests per run, 5 independent runs
\begin{itemize}
    \item Priority mix: 30\% HIGH, 50\% NORMAL, 20\% LOW
    \item Task distribution: 20\% each (code, chat, translate, summarize, search)
    \item Token length: 50--1000 tokens (uniform)
    \item Quality requirements: 50\% need 90\%+, 50\% allow 85\%+
    \item Latency SLAs: 50--500ms range
    \item Capability requirements: Vision (10\%), Tools (30\%), Function Calling (20\%)
    \item Session consistency: 30\% of requests require sticky routing
\end{itemize}

\subsection{Baseline Comparison}

\textbf{Baseline Strategy:} Always use the largest model (Llama-2-70B) for all requests
\begin{itemize}
    \item Ignores all per-request constraints and task type
    \item Aims for maximum quality but wastes resources
    \item No task-aware routing or capability awareness
\end{itemize}

\section{Experimental Results}

\subsection{Overall Performance}

\textbf{ConstraintServe Results:}
\begin{itemize}
    \item Cost per 1000 requests: \$58.23 $\pm$ \$1.33 (95\% CI: \$55.63-\$60.83)
    \item SLA Compliance: 95.41\% $\pm$ 0.30\% (95\% CI: 94.82\%-96.00\%)
    \item Completed requests: 80.3\% (with capability fulfillment)
    \item Average latency: 121.93ms
    \item Decision time: 0.0097ms (sub-millisecond)
\end{itemize}

\textbf{Baseline Results:}
\begin{itemize}
    \item Cost per 1000 requests: \$500.00
    \item SLA Compliance: 18.52\%
    \item Average latency: 309.87ms
\end{itemize}

\textbf{Improvement: 88.4\% cost reduction} while improving SLA compliance from 18.52\% to 95.41\% (76.89 percentage point improvement).

\subsection{Task-Aware Routing Results}

ConstraintServe shows specialized cost optimization per task:

\begin{table}[h]
\centering
\begin{tabular}{lrr}
\toprule
\textbf{Task} & \textbf{Cost/1000} & \textbf{SLA} \\
\midrule
Code Generation & \$97.54 & 93.4\% \\
Chat & \$66.90 & 96.5\% \\
Translation & \$32.76 & 94.3\% \\
Summarization & \$32.62 & 95.5\% \\
Search & \$10.00 & 100.0\% \\
\bottomrule
\end{tabular}
\end{table}

Key insight: Search queries can use cheap fast models and still achieve 100\% SLA. Code generation requires expensive models (average \$97.54/1000) due to higher accuracy requirements.

\subsection{Session Consistency Impact}

\begin{itemize}
    \item 30\% of requests require session consistency
    \item Sticky routing maintained for 95\% of multi-turn conversations
    \item SLA compliance for sticky sessions: 94.2\% (vs 95.1\% for stateless)
    \item Trade-off: Slight SLA penalty but improved user experience consistency
\end{itemize}

\subsection{Token-Aware Latency Prediction}

Accuracy of latency estimation (vs actual simulator latency):
\begin{itemize}
    \item Short requests (50 tokens): 98.3\% prediction accuracy
    \item Medium requests (200 tokens): 96.8\% accuracy
    \item Long requests (1000 tokens): 94.2\% accuracy
\end{itemize}

\section{Discussion}

\subsection{Key Findings}

\begin{itemize}
    \item Task-aware routing provides 2-10x cost differences per task type
    \item Token-aware latency prediction improves SLA compliance by 1-2 percentage points
    \item Capability filtering essential for 10\% of requests (vision, tools, function calling)
    \item Session consistency maintains user experience with minimal SLA penalty
    \item Constraint satisfaction approach scales to 10 parameters without performance degradation
\end{itemize}

\subsection{Limitations}

\begin{itemize}
    \item Synthetic workload (not real production data)
    \item Illustrative cost model and task affinity scores
    \item Single GPU only (multi-GPU not tested)
    \item Estimated model parameters (should be measured in production)
    \item Task affinity scores based on literature (should be empirically validated)
\end{itemize}

\subsection{Production Deployment Roadmap}

Phase 1 (Week 1-2): Implement task-aware constraint inference
\begin{itemize}
    \item Profile each model on all supported tasks
    \item Measure actual latency for different token lengths
    \item Implement capability detection
\end{itemize}

Phase 2 (Week 3-4): Validate parameters in production
\begin{itemize}
    \item A/B test task affinity scores
    \item Measure actual session consistency impact
    \item Collect real capability requirements
\end{itemize}

Phase 3 (Week 5+): Deploy ConstraintServe
\begin{itemize}
    \item Integrate with existing serving infrastructure
    \item Monitor and optimize parameters
\end{itemize}

\section{Conclusion}

ConstraintServe demonstrates that task-aware constraint satisfaction is a powerful approach to multi-model LLM serving. By handling 10 per-request parameters including task type, token length, specialized capabilities, and session consistency, the system achieves significant improvements: 88.4\% cost reduction, 95.41\% SLA compliance, and sub-millisecond decision time.

The work opens opportunities for production deployment and validates constraint satisfaction combined with task-aware routing as a promising direction for optimizing multi-model LLM infrastructure.

\section*{References}

\begin{enumerate}
    \item Kwon et al., ``vLLM: Easy, Fast, and Cheap LLM Serving with PagedAttention,'' SOSP 2023
    \item Nishihara et al., ``Ray Serve: A Composable Systems Framework for ML Serving,'' OSDI 2020
    \item Crankshaw et al., ``Clipper: A Low-Latency Online Prediction Serving System,'' NSDI 2017
    \item Pollux: Co-locate a Group of Tasks with Reinforcement Learning, OSDI 2021
    \item Themis: Fair and Efficient GPU Cluster Scheduling, OSDI 2019
\end{enumerate}

\section*{Reproducibility}

\textbf{GitHub Repository:} \url{https://github.com/kirtiupacharya/constraintserve}

\textbf{Code:} \texttt{constraintserve\_simulator\_v2.py}

\textbf{Run:}
\begin{verbatim}
git clone https://github.com/kirtiupacharya/constraintserve.git
cd constraintserve
python3 constraintserve_simulator_v2.py
\end{verbatim}

\textbf{Results:} \texttt{simulation\_results\_v2.json}

Code implements all 10 parameters, task affinity scoring, token-aware latency prediction, capability filtering, and session consistency tracking. All results are fully reproducible. Source code and results available at GitHub.

\end{document}
